\guidechapter{What a transform actually changes}{Replace a list of values by weights of simple patterns.}

A transform is a change in description. Consider the four numbers
$[8,4,2,2]$. Instead of listing those values directly, describe them as a
sum of four simple patterns. The figure displays each contribution and the
result. Add the first point from each of the first four plots: $4+3+0+1=8$.
At the second point, $4+0+1-1=4$. Repeat at the remaining positions to get
two and two.

\explfigure{15-frequency-patterns}{The same four numbers in a pattern description.}{
Average $[4,4,4,4]$, cosine contribution $[3,0,-3,0]$, sine contribution
$[0,1,0,-1]$, and alternating contribution $[1,-1,1,-1]$ sum exactly to
$[8,4,2,2]$. Lines connect four discrete samples; this is not measured sound.}

Why call these frequencies? The average does not oscillate. The cosine and
sine patterns complete one cycle over the four positions, with different
starting phases. The alternating pattern changes sign at every step.
In this real-pattern basis, the weights are $4,3,1,1$. These are not the
four complex Fourier coefficients in the paper's convention; they are an
equivalent elementary way to display the same real data.

\paragraph{Where complex numbers come in.}
A complex number is a pair of real coordinates written $a+ib$.
Addition works coordinate by coordinate. The rule $i^2=-1$ makes
multiplication by $i$ turn the coordinate plane through a quarter-circle:
$(a,b)$ becomes $(-b,a)$. This is useful bookkeeping for oscillating patterns;
it does not mean the input integers or their final product are imaginary.

\explfigure{24-complex-turns}{A phase is a position around a circle.}{
Multiplication by $i$ cycles through $1,i,-1,-i,1$. Each point has magnitude
one. The four turns provide the simplest visible example of roots of unity.}

For four samples, the ordinary unnormalized Fourier coefficients of
$[8,4,2,2]$ are $[16,6-2i,4,6+2i]$ with the negative-exponent convention.
Dividing by four gives the normalized version
$[4,1.5-0.5i,1,1.5+0.5i]$. The average is four; the paired complex entries
encode the cosine and sine contributions together. You do not need to derive
these complex entries to follow the rest of the guide, but the distinction
explains why transform scale factors cannot be guessed.

\paragraph{Why this helps multiplication.}
Shifting a list changes the phase of each frequency in a simple way.
Convolution is built from weighted shifts, so in the frequency description
it becomes multiplication of matching components, with the chosen normalization
accounted for. The hard work is changing descriptions efficiently and reversing
the change accurately. The full paper adds carefully chosen layouts and a
polynomial representation to address those costs.

\paragraph{What a tensor product means here.}
First apply a two-entry operation to every horizontal pair in a $2\times2$
array, then to every vertical pair. Applying the same operation separately on
both coordinate directions is its tensor-product action. The terminology
becomes compact when there are many axes, but its meaning is still ``do this
operation in each selected direction while holding the others fixed.''

\takeaway{A transform changes the coordinates used to describe data.
Structured patterns can make an operation cheap in those coordinates.}
